% 08_great_research_paper.tex — How to Write a Great Research Paper by Professor Simon Peyton Jones

\section{How to Write a Great Research Paper by Professor Simon Peyton Jones}
\label{sec:great_research_paper}

\subsection*{Writing as a Core Method of Research}

An essential, yet often underestimated, component of research is writing itself.
Contrary to the common view that writing merely records results after the research is done, Dr.\ Simon Peyton Jones emphasizes that writing is central to the research process: ``Writing papers is not just the way that you report the research which you have done, it is how you do the research itself.''
He challenges the typical sequencing of idea $\rightarrow$ do research $\rightarrow$ paper, and instead advocates for idea $\rightarrow$ write paper $\rightarrow$ do research.
Early drafting helps clarify vague ideas, exposes gaps in logic, and facilitates more productive feedback.
Writing becomes a tool for discovering what the research actually is.
``There is no one right way to do research,'' he notes, and writing allows that ambiguity to be captured and resolved.

\subsection*{Clarifying the Central Idea}

A successful research paper revolves around a single central idea that can be clearly communicated and remembered.
``The paper is like a virus,'' Peyton Jones explains.
``You want it to leave your own brain and reach those of others.
You want to convey not a system or an artifact, but an idea.''
Papers that attempt to bundle multiple contributions often dilute their impact; instead, new ideas that emerge during writing can and should be split into future papers.
He offers the example of Mozart---centuries later, people still go to concert halls to ``read his papers.''
This underscores the long-term value of a clear and memorable contribution.
Importantly, he notes that one need not begin with a brilliant idea; even seemingly trivial insights can be refined and elevated through writing.

\subsection*{Structuring the Narrative}

Good research papers ``tell a story.''
This story-like structure should mirror a whiteboard explanation: begin by stating the problem and why it matters, then introduce the core idea, followed by supporting details.
The reader must be convinced that the problem is interesting and that the proposed solution is elegant.
Dr.\ Peyton Jones calls the introduction ``the precious, gold dust page one,'' underscoring its importance for hooking the reader.
He recommends vivid but still formal language, and advises that authors err on the side of readability rather than over-formality.
His reader model---Title (1000 readers), Abstract (100), Introduction (100), Problem (10), Idea (10), Details (3), Related Work (10), and Conclusion (everyone again)---reminds authors to prioritize the earlier sections for maximal clarity and persuasion.

\subsection*{Explicitly Stating Contributions}

To orient the reader and establish expectations, a paper's contributions should be explicitly and concisely listed, ideally on page one.
According to Peyton Jones, a good abstract and introduction should clearly answer two questions: (a) What is the problem you are trying to solve? and (b) What is your contribution to it?
``Molehills not mountains,'' he warns---claims should be modest and defensible.
Each contribution should be falsifiable and directly supported by evidence in the body of the paper.
``Think of stating the contributions as a specification,'' he says, ``and the rest of the paper as the implementation.''
Forward references to relevant sections help readers navigate.
A good litmus test is whether the introduction references all major sections of the paper---if not, either those sections should be cut or they need stronger framing.

\subsection*{Placing Related Work Appropriately}

The related work section should not precede the explanation of the central idea.
When it comes too early, it can confuse or bore the reader---especially when it involves dense citations and domain-specific terminology.
Peyton Jones prefers to place related work near the end of the paper, just before the conclusion, after the reader has enough context to evaluate the novelty and significance of the contribution.
An audience member in the presentation emphasized the importance of a critical comparison with previous work and cautioned against shallow citations.
Peyton Jones further warns against the fallacy of disparaging others to elevate one's own results.
Instead, authors should candidly acknowledge where prior work excels and where their own work may have limitations.
``It is better to state this outright as opposed to leaving a referee or reader to point out this distance.''

\subsection*{Writing for the Reader}

Effective academic writing requires empathy for the reader's experience.
Authors must anticipate background knowledge, preempt confusion, and explain abstract concepts through motivating examples.
Peyton Jones emphasizes that jargon-heavy exposition on page one is a mistake.
Use figures and diagrams liberally to support comprehension.
The goal is not to recount the author's research journey in full, but to communicate the key insights clearly and engagingly.
``Don't drag your readers through the mud,'' he cautions.
The prose should be active, vivid, and focused---avoiding over-technicality while maintaining a level of scholarly rigor.
Readers are not adversaries to be overwhelmed, but collaborators to be guided.

\subsection*{Iterative Feedback and Revision}

Revising based on feedback is indispensable.
Peyton Jones advises sharing drafts early and with diverse readers.
Ask them not just for grammar corrections, but where the logic breaks down---e.g., ``I lost the argument in paragraph 3.''
Peers unfamiliar with the topic can provide fresh perspective; experts can offer precise critiques, especially when quoted or compared.
``Try to make it so that they could not possibly not misunderstand,'' he urges.
When receiving peer reviews---especially negative ones---he recommends waiting a few days before reacting, allowing emotions to settle.
Misunderstandings often stem from unclear writing, not unqualified reviewers.
Lastly, when describing the work of others, he suggests sending the draft to those authors to check for fairness---collaboration over competition fosters better scholarship.
